\documentclass[../DoAn.tex]{subfiles}
\begin{document}

\section{Xây dựng ứng dụng}

\subsection{Thư viện và công cụ sử dụng}
Quá trình xây dựng ứng dụng Ví Stealth Address được phát triển trên nhiều nền tảng và ngôn ngữ khác nhau, từ xử lý logic blockchain đến giao diện web và các thuật toán toán học Zero-Knowledge. Danh sách các thư viện và công cụ được liệt kê chi tiết trong Bảng \ref{table:tools}.

\begin{table}[H]
\centering{}
\resizebox{\textwidth}{!}{
    \begin{tabular}{|l|l|l|l|}
        \hline
        \textbf{Mục đích} & \textbf{Công cụ} & \textbf{Phiên bản} & \textbf{Địa chỉ URL} \\ \hline
        Ngôn ngữ lập trình Smart Contract & Solidity & v0.8.28 & https://soliditylang.org/ \\ \hline
        Khung làm việc (Framework) EVM & Hardhat & v2.22.x & https://hardhat.org/ \\ \hline
        Ngôn ngữ viết mạch ZK-SNARK & Circom & v2.0.0 & https://docs.circom.io/ \\ \hline
        Thư viện sinh và xác minh ZK Proof & SnarkJS & v0.7.0 & https://github.com/iden3/snarkjs \\ \hline
        Môi trường chạy Server Relayer & Node.js (Express) & v20.x & https://nodejs.org/ \\ \hline
        Thư viện giao tiếp Web3 / Mật mã & ethers.js & v6.11.1 & https://docs.ethers.org/v6/ \\ \hline
        Phát triển giao diện Frontend & React (Vite) & v18.2.0 & https://react.dev/ \\ \hline
    \end{tabular}
}
    \caption{Danh sách thư viện và công cụ sử dụng}
    \label{table:tools}
\end{table}

\subsection{Kết quả đạt được}
Kết quả thực nghiệm đã đóng gói thành công 4 phân hệ (module) chính của hệ thống, tương tác chặt chẽ với nhau theo mô hình Account Abstraction:
\begin{itemize}
    \item \textbf{Gói Blockchain (\texttt{blockchain/contracts}):} Đóng gói các hợp đồng thông minh lõi, bao gồm \texttt{StealthAccount.sol} (Ví AA), \texttt{ERC5564Announcer.sol} (Lưu vết giao dịch), \texttt{StealthTreeManager.sol} (Quản lý Merkle Tree on-chain) và \texttt{SocialRecoveryModule.sol} (Quản lý khôi phục xã hội).
    \item \textbf{Gói Frontend (\texttt{stealth-wallet}):} Giao diện người dùng Web Application, đảm nhiệm việc sinh khóa cục bộ (Local Key Generation), tự động quét các sự kiện trên chuỗi và tích hợp WebAssembly (WASM) để sinh bằng chứng ZK ngay trên trình duyệt mà không làm rò rỉ khóa riêng.
    \item \textbf{Gói Mạch ZK (\texttt{zk/circuits}):} Bao gồm các mạch \texttt{stealth.circom} và \texttt{smt\_update.circom} được biên dịch thành file WASM và zkey, phục vụ cho việc sinh và xác minh bằng chứng mật mã.
    \item \textbf{Gói Relayer Backend (\texttt{server}):} Server API viết bằng Node.js xử lý việc nhận các \texttt{UserOperation}, duy trì cơ sở dữ liệu cây Merkle off-chain (\texttt{leaves.json}) và thay mặt người dùng trả phí gas.
\end{itemize}

Thông tin thống kê về quy mô mã nguồn của ứng dụng được thể hiện qua Bảng \ref{table:stats}.

\begin{table}[H]
\centering{}
    \begin{tabular}{|l|c|c|c|}
        \hline
        \textbf{Phân hệ / Sản phẩm đóng gói} & \textbf{Số lượng tệp tin} & \textbf{Số dòng code (LOC)} & \textbf{Dung lượng mã nguồn} \\ \hline
        Smart Contracts (Solidity) & 12 & $\sim$ 1,500 & 85 KB \\ \hline
        Mạch ZK (Circom) & 5 & $\sim$ 500 & 20 KB \\ \hline
        Frontend Web App (React/JS) & 25 & $\sim$ 2,500 & 2.5 MB \\ \hline
        Relayer Server (Node.js) & 10 & $\sim$ 800 & 150 KB \\ \hline
        \textbf{Tổng cộng (Ước tính)} & \textbf{52} & \textbf{$\sim$ 5,300} & \textbf{$\sim$ 2.75 MB} \\ \hline
    \end{tabular}
    \caption{Thống kê quy mô mã nguồn ứng dụng}
    \label{table:stats}
\end{table}

\subsection{Minh họa các chức năng chính}

Phần này trình bày giao diện của các tính năng quan trọng nhất trong hệ thống kèm theo mô tả về luồng dữ liệu bên dưới.

\begin{figure}[H]
    \centering
    \includegraphics[width=0.8\textwidth]{Hinhve/placeholder_create_wallet.png}
    \caption{Giao diện khởi tạo Stealth Meta-Address}
    \label{fig:func1}
\end{figure}
\textbf{Giải thích:} Tại giao diện này, hệ thống sẽ sử dụng thuật toán đường cong elliptic \texttt{secp256k1} qua thư viện \texttt{ethers.js} để tạo hai cặp khóa: Khóa xem (Viewing Key) và Khóa chi tiêu (Spending Key). Hệ thống nối hai khóa công khai lại và mã hóa Base58 để hiển thị một \texttt{Stealth Meta-Address} duy nhất cho người dùng. Khóa cá nhân được mã hóa và lưu an toàn ở bộ nhớ cục bộ (Local Storage).

\begin{figure}[H]
    \centering
    \includegraphics[width=0.8\textwidth]{Hinhve/placeholder_send_transfer.png}
    \caption{Giao diện chuyển tiền ẩn danh (Stealth Transfer)}
    \label{fig:func2}
\end{figure}
\textbf{Giải thích:} Người gửi nhập Meta-Address của người nhận. Dưới nền tảng, ứng dụng tự động sinh một khóa tạm thời (Ephemeral Key), kết hợp với Khóa xem công khai của người nhận để sinh ra Bí mật chung (Shared Secret) theo giao thức ECDH. Địa chỉ ẩn danh cuối cùng được tạo ra và giao dịch được gửi lên chuỗi, đính kèm sự kiện \texttt{Announcement} ghi nhận Ephemeral Public Key.

\begin{figure}[H]
    \centering
    \includegraphics[width=0.8\textwidth]{Hinhve/placeholder_scanner.png}
    \caption{Giao diện quét và khám phá tài sản (Scanner)}
    \label{fig:func3}
\end{figure}
\textbf{Giải thích:} Giao diện cho phép người dùng khởi chạy tiến trình quét (sử dụng Web Worker để không gây giật lag trình duyệt). Trình quét tải toàn bộ sự kiện \texttt{Announcement} từ blockchain, kết hợp với Khóa xem bí mật để đối chiếu View Tag (1-byte). Nếu khớp, hệ thống giải mã để hiển thị số dư thực tế đang nằm trong các địa chỉ ẩn danh thuộc về người dùng.

\begin{figure}[H]
    \centering
    \includegraphics[width=0.8\textwidth]{Hinhve/placeholder_zk_spend.png}
    \caption{Giao diện rút tiền ẩn danh (ZK Spend)}
    \label{fig:func4}
\end{figure}
\textbf{Giải thích:} Khi người dùng yêu cầu rút tiền từ ví ẩn danh về ví sàn giao dịch, ứng dụng sẽ thu thập Merkle Root từ server, tính toán đường dẫn Merkle Path và chạy thư viện \texttt{snarkjs} (dưới dạng WASM) để sinh bằng chứng toán học. Toàn bộ quá trình tạo Proof chỉ tốn vài giây. Proof sau đó được gửi off-chain cho Relayer, Relayer sẽ trả phí gas và gửi giao dịch \texttt{execute} lên mạng lưới.

\section{Kiểm thử}
Quá trình kiểm thử tập trung vào hai nghiệp vụ lõi phức tạp nhất của hệ thống: Chi tiêu ZK (Stealth Spend) và Khôi phục mạng lưới xã hội (Social Recovery). Hệ thống sử dụng thư viện \texttt{Chai} và \texttt{Hardhat Testing} để viết unit test on-chain. Tổng cộng có 25 kịch bản kiểm thử (Test Cases) đã được chạy với tỷ lệ thành công (Pass rate) đạt 100\%.

Dưới đây là một số trường hợp kiểm thử (Test Cases) tiêu biểu.

\begin{table}[H]
\centering{}
\resizebox{\textwidth}{!}{
    \begin{tabular}{|p{1cm}|p{2cm}|p{4cm}|p{4cm}|p{3cm}|p{1.5cm}|}
        \hline
        \textbf{Mã TC} & \textbf{Chức năng} & \textbf{Các bước thực hiện} & \textbf{Kết quả mong đợi} & \textbf{Kết quả thực tế} & \textbf{Trạng thái} \\ \hline
        TC\_01 & ZK Spend & 1. Nhập đúng \texttt{spendPriv}. 2. Sinh ZK Proof. 3. Giao dịch \texttt{execute} tới Smart Contract. & \texttt{verifier.verifyProof} trả về \texttt{true}. Giao dịch chuyển ETH thành công. & Chuyển ETH thành công. & Pass \\ \hline
        TC\_02 & ZK Spend & 1. Nhập \texttt{spendPriv} sai (của người khác). 2. Sinh ZK Proof. 3. Giao dịch \texttt{execute}. & Proof không hợp lệ. Smart Contract revert lỗi \texttt{INVALID\_PROOF}. & Bị revert đúng lỗi chỉ định. & Pass \\ \hline
        TC\_03 & ZK Spend (Replay) & 1. Bắt gói tin ZK Proof cũ. 2. Đổi \texttt{indexCommitment} sang ví khác và gửi lại. & ZK Proof bị từ chối do không khớp tín hiệu Public Inputs gốc. & Bị revert lỗi tín hiệu không khớp. & Pass \\ \hline
        TC\_04 & Social Recovery & 1. 1 Guardian gọi \texttt{proposeRecovery}. 2. 2 Guardians gọi \texttt{approveRecovery} (Threshold=3). 3. Bất kỳ ai gọi \texttt{executeRecovery} kèm ZK Proof cập nhật Merkle. & Root mới được cập nhật trên \texttt{StealthTreeManager}. Ví thay đổi Owner thành công. & Cập nhật Root và Leaf thành công. & Pass \\ \hline
        TC\_05 & Social Recovery & 1. 1 Guardian gọi \texttt{proposeRecovery}. 2. 1 Guardian gọi \texttt{approveRecovery}. 3. Gọi \texttt{executeRecovery}. & Bị revert do số lượng phê duyệt (2) nhỏ hơn Threshold (3). & Bị revert lỗi \texttt{not enough approvals}. & Pass \\ \hline
    \end{tabular}
}
    \caption{Các trường hợp kiểm thử nghiệp vụ lõi (Test Cases)}
    \label{table:testcases}
\end{table}

\textbf{Tổng kết kiểm thử:} Quá trình kiểm tra đảm bảo hệ thống không tồn tại lỗ hổng rò rỉ khóa trong quá trình sinh bằng chứng, và bảo vệ tuyệt đối khỏi các dạng tấn công lặp lại (Replay Attacks) cũng như tấn công thay đổi trạng thái (State Manipulation) từ người bảo hộ giả mạo.

\section{Triển khai}
\subsection{Mô hình triển khai thử nghiệm}
Môi trường thực nghiệm được thiết lập chạy cục bộ (Localhost) để giả lập mạng lưới, nhằm phục vụ công tác tinh chỉnh và đo đạc thông số:
\begin{itemize}
    \item \textbf{Mạng Blockchain:} Triển khai trên node cục bộ \texttt{Hardhat Network} (ChainID: 31337). Quá trình deploy sinh ra 5 hợp đồng thông minh lõi và hệ thống tự động ghi địa chỉ hợp đồng vào biến môi trường của Frontend.
    \item \textbf{Server Backend:} Chạy trên nền tảng Node.js (cấu hình thử nghiệm giả lập: 4 Cores CPU, 8GB RAM). Server mở port 3001, lắng nghe yêu cầu khởi tạo cây Merkle và trả kết quả API dưới dạng JSON.
    \item \textbf{Trình duyệt Client:} Giao diện Web chạy qua Vite (port 5173). Thử nghiệm thực hiện trên trình duyệt Google Chrome v120.
\end{itemize}

\subsection{Kết quả và Thông số kỹ thuật đạt được}
Trong quá trình triển khai thử nghiệm, hệ thống đã ghi nhận các thông số hiệu năng (Performance Metrics) ở mức xuất sắc đối với một ứng dụng Web3 mật mã hóa cao:

\textbf{1. Hiệu năng tính toán ZK Proof ở Client:}
\begin{itemize}
    \item Độ sâu cây Merkle (Levels): 20 (Hỗ trợ hơn 1 triệu địa chỉ ẩn danh).
    \item Thời gian tải thư viện \texttt{snarkjs} và file WASM (2MB): $\sim 0.5$ giây.
    \item \textbf{Thời gian sinh bằng chứng (Generate Proof):} Khoảng $\sim 1.5 - 2.5$ giây. Đây là kết quả rất tốt, đảm bảo trải nghiệm người dùng (UX) không bị gián đoạn quá lâu.
\end{itemize}

\textbf{2. Chi phí Gas (Phí giao dịch On-chain):}
Mọi tính toán ZK được xác minh qua các Precompiled Contracts tích hợp sẵn trên EVM (như phép toán trên đường cong alt\_bn128/bn254).
\begin{itemize}
    \item Giao dịch \texttt{execute} (Xác minh ZK Spend + Chuyển tiền): $\sim 310,000$ Gas.
    \item Đề xuất khôi phục \texttt{proposeRecovery}: $\sim 95,000$ Gas.
    \item Thực thi khôi phục \texttt{executeRecovery} (Xác minh cập nhật SMT Merkle): $\sim 350,000$ Gas.
    \item \textit{Đánh giá:} Chi phí này rất cạnh tranh và tương đương với các giải pháp máy trộn (Mixer) nổi tiếng trên thị trường, nhưng mang lại mức độ phân quyền và bảo vệ định tuyến an toàn hơn.
\end{itemize}

\textbf{3. Tốc độ Quét (Scanning Speed):}
Nhờ kiến trúc \textit{View Tag Optimization} (chỉ đối chiếu byte đầu tiên của mã băm thay vì nhân ma trận toàn khối lượng giao dịch), tốc độ quét mạng lưới tăng gấp 255 lần.
\begin{itemize}
    \item Quét thử nghiệm 1,000 sự kiện \texttt{Announcement}: Xử lý cục bộ mất chưa đến 0.2 giây.
    \item Quét thử nghiệm 10,000 sự kiện: Chỉ mất $\sim 1.8$ giây, chứng minh hệ thống có khả năng chịu tải dữ liệu lịch sử cực lớn khi triển khai thực tế trên mạng Mainnet.
\end{itemize}

\end{document}
